% GDZ PPN236018647,native526–527;HKUST527–528.
\clearpage\SourceLeaf{526}\ArithmeticHead{497}{R}{1796 IUN. 22.—1796 IUL.}
\par\medskip\Needspace{4\baselineskip}\begin{center}[18.\footnotemark]\end{center}
\begingroup\renewcommand{\thefootnote}{[*)}\footnotetext{Zwischen den Notizen Nr. 17 und Nr. 18 steht in der Handschrift noch eine Aufzeichnung, die \textsc{Gauss} durchstrichen hat und die darum nur schwer lesbar ist (vergl. die Nachbildung). Sie lautet wie folgt:]
\begin{center}[17a.]\end{center}
\indent Summa trium quadratorum continue proportionalium numquam primus esse potest: conspicuum exemplum novimus et quod congruum videtur. Confidamus.
\par\noindent\hfill[1796] 9. Iul.
\par\medskip\indent Wenn die drei Quadrate ganzer Zahlen \(x^2,y^2,z^2\) in laufender Proportion stehen, ist
\[
x:y:z=1:\frac{m}{n}:\frac{m^2}{n^2}.
\]
\noindent Als Summe der drei Quadrate ergibt sich also die zerlegbare Form
\[
m^4+m^2n^2+n^4=(m^2+mn+n^2)(m^2-mn+n^2),
\]
\noindent die \textls[120]{im allgemeinen} keine Primzahl darstellt. Vielleicht hat \textsc{Gauss} zunächst die Möglichkeit
\[
m^2+mn+n^2=p,\quad m^2-mn+n^2=1
\]
\noindent übersehen, die für \(m=n=1,\ p=3\) eintritt. Bei dem »conspicuum exemplum« war wohl für \(m\) oder \(n\) eine von Eins verschiedene ganze Zahl genommen worden.
\par\noindent\hfill\textsc{Bachmann.}}
\endgroup
\indent {\fontspec{Times New Roman}ΕΥΡΗΚΑ!} num[erus]\(=\triangle+\triangle+\triangle\).
\par\noindent\hfill[1796] 10. Iul.\quad Gott[ingae]

\medskip\begingroup\small
\indent Vergl. die Bemerkung zu der Nr. 17; es soll zum Ausdruck gebracht werden, daß jede Zahl als Summe von drei Dreieckszahlen dargestellt werden kann.
\par\noindent\hfill\textsc{Klein. Bachmann.}
\par\endgroup

\NumberedItem{19}
\indent Determinatio Euleriana formarum in quibus numeri compositi plus una vice continentur.
\par\noindent\hfill[1796 Iul.\quad Gottingae]

\medskip\begingroup\small
\indent Gemeint sind hier jedenfalls die Methoden, die \textsc{Euler} für die Entscheidung der Frage angegeben hat, ob eine vorgegebene Zahl von der Form \(4n+1\) eine Primzahl ist oder nicht; siehe die §§ 40—43 der Abhandlung \textit{De numeris, qui sunt aggregata duorum quadratorum}, Novi Comm. Acad. Petropol. \textbf{4} (1752/53) 1758, S. 3—40, insbesondere S. 29 ff.\begingroup\renewcommand{\thefootnote}{**)}\footnote{L. Euleri Opera omnia, series I, vol. 2, S 295. Vergl. auch noch die Abhandlung \textit{Quomodo numeri praemagni sint explorandi, utrum sint primi nec ne}, Novi Comm. Acad. Petrop. \textbf{13} (1768) 1769, S. 67—88, Opera omnia, ser. I, vol. 3, S. 112.}\endgroup\ und \textit{Extrait d’une lettre de M. Fuss à M. Béguelin}, Nouv. Mémoires de l’Académie de Berlin (1776) 1779, S. 340—346\begingroup\renewcommand{\thefootnote}{***)}\footnote{L. Euleri Opera omnia, ser. I, vol. 3, S. 421; diese Note ist ein ausführlicher Bericht über den Inhalt der erst nach \textsc{Eulers} Tode veröffentlichten, 1801 erschienenen Abhandlung \textit{De formulis speciei \(mxx+nyy\) ad numeros primos explorandos idoneis, earumque mirabilibus proprietatibus}, Nova Acta Acad. Petrop. \textbf{12} (1794) 1801, S. 22—46, die also \textsc{Gauss} im Juli 1796 nicht gekannt haben kann.}\endgroup. Im \textsc{Leiste} finden sich an mehreren Stellen nach diesen Methoden ausgeführte Rechnungen, so heißt es z. B. bei S. 48:
\par\endgroup
\par\smallskip\noindent\small X 1.\hfill 63\normalsize

\clearpage\SourceLeaf{527}\ArithmeticHead{498}{L}{TAGEBUCH MIT ERLÄUTERUNGEN.}
\begingroup\small
\indent Discerpend. 283009 in bina quadrata
\[
4225+528^2,
\]
\noindent also Primzahl.
\par\noindent\hfill\textsc{Stäckel. Schlesinger.}
\par\endgroup

\NumberedItem{20}
\indent Principia componendi scalas serierum variatim recurrentium.
\par\noindent\hfill[1796] 16. Iul.\quad Gott[ingae]

\medskip\begingroup\small
\indent Vergl. die Bemerkung zu der Nr. 8. Daß \textsc{Gauss} sich auch noch im folgenden Monat mit rekurrenten Reihen beschäftigt hat, zeigt der vom 28. August 1796 datierte art. [3.] der \textit{Exercitationes mathem.}, oben S. 139.
\par\noindent\hfill\textsc{Loewy.}
\par\endgroup

\NumberedItem{21}
\indent Methodus Euleriana pro demonstranda relatione inter rectangula sub segmentis rectarum sese secantium in sectionibus conicis ad omnes curvas applicata[\footnotemark].
\begingroup\renewcommand{\thefootnote}{[*)}\footnotetext{In der Handschrift steht \textls[120]{applicatum}.]}\endgroup
\par\noindent\hfill[1796] 31. Iul.\quad Gott[ingae]

\medskip\begingroup\small
\indent \textsc{Euler} gibt in der \textit{Introductio in analysin infinitorum}, Lausannae 1748, II, § 92, 93 einen einfachen Beweis für den im wesentlichen schon von \textsc{Apollonius} (Buch 3, § 17, 19, 22) erkannten Satz: Werden durch einen Punkt \(O\) zwei Geraden gezogen, die einen Kegelschnitt in den Punkten \(A,\ B\) und \(A',\ B'\) treffen, so hat, wenn nur die Geraden ihre Richtung beibehalten, das Verhältnis der Produkte \(OA.OB\) und \(OA'.OB'\) für alle Lagen von \(O\) denselben Wert. Daß der Beweis sich auf Kurven höherer Ordnung übertragen läßt, ist \textsc{Euler} nicht entgangen; an einer spätern Stelle der \textit{Introductio} hat er dies für die Kurven dritter Ordnung ausführlich dargelegt und sagt dann (§ 247): Atque similis proprietas in lineas quarti, quinti atque superiorum ordinum competet.
\par\noindent\hfill\textsc{Stäckel.}
\par\endgroup

\NumberedItem{22}
\indent \(a^{2^n\mp1\,(p)}\equiv1\) semper solvere in potestate.
\par\noindent\hfill[1796] Aug. 3.\quad Gott[ingae]

\medskip\begingroup\small
\indent Diese Aufzeichnung besagt vermutlich, daß \textsc{Gauss}, wenn \(p\) eine Primzahl von der Form \(2^n\mp1\) ist, die Funktion \(x^p-1\) in Bezug auf einen Primzahlmodul in ihre Primfunktionfaktoren zu zerlegen vermochte. Vergl. \textit{Analysis Residuorum}, Werke II, S. 229 und S. 199—211; S. 209 wird der Fall \(p=31=2^5-1\) modulo 331 behandelt.
\par\noindent\hfill\textsc{Bachmann.}
\par\endgroup
